\documentclass[conference]{IEEEtran}
\IEEEoverridecommandlockouts
\usepackage{amsmath,amssymb}
\usepackage{booktabs}
\usepackage{placeins}
\usepackage{cite}
\usepackage{tikz}
\usetikzlibrary{arrows.meta}
\usepackage{url}
\usepackage{xeCJK}
\setmainfont{TeX Gyre Termes}
\setCJKmainfont{FandolSong}
\setCJKsansfont{FandolHei}
\usepackage[hidelinks]{hyperref}

\title{TEMPO：用于市场中性对冲的时间证据与动量}
\author{\IEEEauthorblockN{Yicheng He\textsuperscript{1,2} and Qianqian Xie\textsuperscript{1,2,*}}
\IEEEauthorblockA{\textsuperscript{1}School of Artificial Intelligence, Wuhan University, Wuhan, China\\
\textsuperscript{2}Cognitive Language and Information Lab, Wuhan University, Wuhan, China\\
Team TEMPO\\
xqq.sincere@gmail.com}
\thanks{\textsuperscript{*}通讯作者。}}

\begin{document}
\maketitle

\begin{abstract}
本文提出市场中性对冲智能体 TEMPO，将时间证据与相对价格动量结合，用于 FinReason Cup 2026 任务二。语言模型解释新闻事件和已经结算的预测误差。数值模块将时间证据表示为具有时间衰减的事件、股票关系和反馈因子，计算这些因子在规定持仓区间内的贡献，并将其预测与已观察到的价格变化结合，选择金额相等的做多、做空股票组合，或者放弃交易。我们使用 2026 年 10 月七个公开发布日中具有进场和退出价格的 37 个轮次进行历史回放。每次完整交易收取账户权益的 10 个基点作为费用时，TEMPO 在 30 个活跃轮次中取得 9.66\% 的累计净收益，按结算时点计算的最大回撤为 1.14\%。仅使用因子执行交易取得 4.67\%。组成部分消融实验中，仅使用价格动量取得 3.15\%，加入事件因子后取得 9.40\%，再加入关系因子后取得 9.66\%；在本次回放中，加入反馈因子没有改变交易行动。将每次完整交易的费用提高到 20 个基点后，保持 TEMPO 的交易行动不变，累计净收益为 6.43\%。在 11 月 2 日独立计分的六轮实时交易中，事前固定的方法取得 2.275\% 的官方收益、1.294 的非年化 Sharpe，以及零结算时点回撤，在参赛队伍中排名第二，收益超过全部三种官方基线。这一单日结果提供了有限的前瞻性证据。
\end{abstract}

\begin{IEEEkeywords}
金融智能体，时间证据，价格动量，市场中性对冲，语言模型
\end{IEEEkeywords}

\section{引言}
FinReason Cup 2026 任务二要求参与者根据已经发布的价格和简短新闻摘要，在固定的股票范围内反复作出决策\cite{finreason2026overview,finreason2026task2overview}。决策产生后，实际持仓开始之前，依据可能已经失去时效：当前规则在下一个估值时点进场，在再下一个估值时点退出\cite{task2protocol}。因此，智能体需要把证据出现的时间与交易接受评价的持仓区间联系起来。HERCULEAN 提供了更广泛的金融智能体评测背景，并将对冲列为一项独立工作任务\cite{peng2026herculeanagenticbenchmarkfinancial}。

本文提出 TEMPO（Temporal Evidence and Momentum for Market-Neutral Hedging），即用于市场中性对冲的时间证据与动量方法。TEMPO 通过一组只追加的带符号因子表示时间证据。每个因子都有创建时间、幅度、衰减类型和时间尺度。它对未来持仓区间的贡献，是该区间内尚未释放的影响所对应的变化量。因子来自三个来源：新闻事件、相关股票之间的条件偏离，以及结算后观察到的预测误差。动量部分表示当日交易时段内最新的相对价格变化，决策模块将其与基于证据的预测结合。语言模型负责分类解释；确定性的数值计算负责产生预测和配对交易行动。

历史实验考察延迟执行条件下，时间证据和价格动量怎样影响交易收益。我们将 TEMPO 与仅使用因子的交易方法、情绪加动量的参考方法比较，逐类移除因子以衡量其对股票选择的贡献，并考察动量权重和交易费用的影响。所有方法使用相同的进场和退出价格接受评价。随后，我们将事前固定的方法用于 11 月 2 日独立计分的实时交易，并将官方六轮成绩与 10 月研究回放分别报告。

\section{相关研究}
价格变化是否反转，与变化的起因和观察期限有关。Engelberg、Gao 和 Jagannathan 将配对交易结果与流动性、共同信息以及个股新闻联系起来；由具有价值意义的个股新闻引起的价格分歧，对价格重新接近的交易较为不利\cite{engelberg2008}。Heston、Korajczyk 和 Sadka 发现，以报价衡量的短期反转主要集中在最初的半小时滞后期，流动性恢复所需时间超过 30 分钟且短于 60 分钟\cite{heston2010}。这些发现支持分别表示新闻事件和股票关系偏离，并在任务规定的未来持仓区间内评价每个信号。

Tetlock 发现，市场对旧信息的反应与对新信息的反应存在差异\cite{tetlock2011}。TEMPO 通过事件创建时间和衰减函数明确表示已经观察到的信息有多旧。任务提供的压缩新闻摘要构成事件模型的文本输入。

条件因子模型将系统性收益与剩余偏离分开\cite{spears2023}。TEMPO 的关系模块测量十只股票两两之间的条件偏离，并使用 Ledoit--Wolf 收缩协方差估计\cite{ledoit2004}，确定回归系数、相关性准入条件和残差波动阈值。

谢倩倩老师及合作者围绕金融语言模型与评测，研究了 PIXIU、FinBen、Open-FinLLMs、多语言金融任务、金融专业能力、多模态输入、推理和结构化信息提取\cite{xie2023pixiu,xie2024finben,xie2024openfinllms,hu2024nolang,peng2025multifinben,yang2025ucfe,lin2025leaderboard,qian2025fino1,cao2025finaudio,kuang2025scores,wang2025fintagging,zhang2024dolares,peng2025plutus}。其中，股票涨跌预测与金融时间序列预测研究，与本文结合文本证据和近期价格的设计密切相关\cite{xie2023wallstreet,xiao2025retrieval}。金融决策代理和投资顾问评测为行动选择提供了更广泛的研究背景\cite{yu2024fincon,li2025investorbench,xiong2025flagtrader,gong2026shijian,yang2026fincua}。配对交易研究涉及股票选择及考虑风险的连续决策\cite{han2023selecttrade,han2023riskpair}。其他金融应用包括信用评分、审计、税务、虚假信息检测、金融模拟以及共享任务设计\cite{feng2023credit,huang2024auditwen,huang2026auditjudg,hu2026taxpraben,liu2025fmdllama,liu2025misinformation,liu2026simulation,xie2024agentScen}。这些工作展示了金融语言和决策研究的范围；本文针对比赛公开轮次评价所述交易规则。

\section{任务、数据与时间访问规则}
\subsection{决策与结算规则}
股票范围包括 NVDA、GOOG、META、MSFT、JPM、AMZN、AAPL、AMD、WMT 和 TSLA。对于轮次 $t$，$p_t$ 表示实际发布时间，$e_t$ 和 $x_t$ 分别表示预定进场时间和退出时间。$P_{i,u}$ 表示股票 $i$ 在时间 $u$ 公布的估值价格。TEMPO 只能读取 $p_t$ 时刻及此前已经发布的内容。发生交易时，账户权益的一半分配给做多股票 $L_t$，另一半分配给做空股票 $S_t$。交易的毛收益率为
\begin{equation}
 R_t^{\mathrm{gross}}=\tfrac12\left(
 \frac{P_{L_t,x_t}}{P_{L_t,e_t}}-
 \frac{P_{S_t,x_t}}{P_{S_t,e_t}}\right).
\end{equation}
本地研究目标采用
\begin{equation}
 R_t^{\mathrm{net}}=R_t^{\mathrm{gross}}-c,
 \qquad c=0.001,
\end{equation}
其中费用仅在发生交易的轮次扣除；空仓轮次的收益和费用均为零。本地回放对一轮完整配对交易收取账户权益的 10 个基点。决策规则还要求预测收益额外超过费用 5 个基点。

所有时间均保留原有的时区偏移。只有退出价格的发布完成以后，相应预测误差才能用于反馈。若缺少最终价格，该持仓区间不会进入收益评价，也不会产生反馈目标。

\subsection{公开发布内容与历史价格}
10 月研究数据快照包含 2026 年 10 月 14 日、15 日、19 日、20 日、21 日、22 日和 23 日的 41 次发布，共有 410 条股票摘要。41 个决策轮次中，37 个轮次的两个规定估值时点都能在公开价格文件中找到。其余四个轮次为第 11、12、40 和 41 轮；这些轮次的最终估值价格不在已经发布的数据快照内。它们可以产生面向未来的交易行动，但不会计入实测收益，也不会成为反馈目标。

提供的新闻内容包括长度为 30--50 个 token 的英文摘要、检索到的文章数量、覆盖范围元数据和来源列表。事件模型接收每只股票的摘要和检索文章数量；元数据用于识别内容覆盖不完整的情况。

数值参数校准使用已经缓存的 Yahoo Finance 小时行情开盘价，覆盖 2023 年 12 月 4 日至 2026 年 10 月 13 日的 709 个完整交易日。每个交易日都包含纽约时间 09:30、10:30、11:30、12:30、14:30 和 15:30 六个共同观测时点。估计协方差之前，将 3,545 个日内对数收益向量除以对应时间间隔小时数的平方根。隔夜收益不参与这一协方差估计。历史小时行情用于确定波动率和协方差参数；回放收益使用比赛公布的估值价格。

回放检查发布编号、时间戳、发布截止时间和可获得文件的哈希值。历史参数校准在比赛首次发布之前结束；反馈只使用当前决策时刻之前已经发布的结果。

\section{TEMPO 架构}
每次发布时，TEMPO 先评价新近结算的预测，并追加模型接受的反馈因子。随后，它根据已经发布的价格和新闻生成关系因子与事件因子。这些因子表示时间证据，并产生预定持仓区间的预测。决策模块将预测与相对价格动量结合，选择股票组合或者放弃交易。图~\ref{fig:architecture} 展示信息流。

\begin{figure*}[t]
\centering
\begin{tikzpicture}[
  box/.style={draw, line width=0.6pt, align=center, minimum height=9mm,
              text width=2.55cm, inner sep=3pt, font=\footnotesize},
  flow/.style={-{Latex[length=2mm]}, line width=0.6pt}
]
\node[box] (news) at (1.5,2.0) {已发布的\\新闻摘要};
\node[box] (prices) at (1.5,0.0) {已发布的\\估值价格};
\node[box] (settled) at (1.5,-2.0) {已公布的结算结果\\与保存的预测};

\node[box] (events) at (4.7,2.0) {语言模型：\\事件因子};
\node[box] (relations) at (4.7,0.55) {股票组合的\\条件偏离};
\node[box] (momentum) at (4.7,-0.9) {相对价格\\动量};
\node[box] (feedback) at (4.7,-2.35) {语言模型：\\反馈因子};

\node[box] (ledger) at (8.1,1.25) {只追加的事件、\\关系与反馈因子};
\node[box] (interval) at (11.45,1.25) {计算 $[e_t,x_t]$ 内的\\衰减并汇总预测\\$g_{i,t},\,\mu_{i,t}$};
\node[box] (decision) at (14.8,0.25) {结合预测与动量；\\排列股票并检查\\收益费用条件};
\node[box] (action) at (14.8,-1.45) {金额相等的做多、\\做空组合或放弃交易};

\draw[flow] (news) -- (events);
\draw[flow] (prices.east) -- (relations.west);
\draw[flow] (prices.east) -- (momentum.west);
\draw[flow] (settled.east) -- (feedback.west);
\draw[flow] (events.east) -- (ledger.west);
\draw[flow] (relations.east) -- (ledger.west);
\draw[flow] (feedback.east) to[out=0,in=-90] (ledger.south);
\draw[flow] (ledger) -- (interval);
\draw[flow] (interval.east) -- (decision.north west);
\draw[flow] (momentum.east) to[out=0,in=180] (decision.west);
\draw[flow] (decision) -- (action);
\end{tikzpicture}
\caption{TEMPO 在每次发布时的信息流。只有相应退出价格公布以后，结算结果才会进入反馈路径；因子的贡献在预定的进场至退出区间内计算。}
\label{fig:architecture}
\end{figure*}

\subsection{时间证据及其在持仓区间内的贡献}
因子 $k$ 保存所属股票 $i(k)$、创建时间 $p_k$、带符号幅度 $A_k$、时间尺度 $\tau_k>0$ 和类型 $z_k$。创建时间是相应发布内容可供智能体使用的时间。因子建立以后，其幅度和方向不会修改。设实际经过的小时数为 $u\geq0$，尚存影响的函数为
\begin{align}
 S_{\mathrm{event}}(u;\tau)&=\exp[-(u/\tau)^{0.75}],\\
 S_{\mathrm{relation}}(u;\tau)&=\exp(-u/\tau),\\
 S_{\mathrm{feedback}}(u;\tau)&=(1+u/\tau)^{-1.5}.
\end{align}
该因子对预测收益的贡献为
\begin{equation}
 v_{k,t}=A_k\left[S_{z_k}(e_t-p_k;\tau_k)
                    -S_{z_k}(x_t-p_k;\tau_k)\right].
 \label{eq:interval}
\end{equation}
式~\ref{eq:interval} 中的时间差以小时表示。每个尚存影响函数都会随时间下降，所以单个因子的贡献保持创建时的方向。该式衡量进场至退出期间尚存影响的变化；因子在隔夜期间也会按照实际时间继续老化。

三种衰减函数表示不同的持续时间假设：拉伸指数函数使新闻事件在早期产生较快反应，同时保留较长的尾部；指数函数为股票关系偏离设置单一特征时间；幂函数允许结算后的预测误差保留较久影响。函数形式及指数是建模选择，37 轮样本无法单独确定它们。三种函数的半衰时间依次为 $\tau(\ln 2)^{4/3}$、$\tau\ln 2$ 和 $\tau(2^{2/3}-1)$。比赛规则已经确定进场和退出时间，因此衰减使用实际经过的时间。

\subsection{新闻事件因子}
每次发布只进行一次语言模型请求，同时处理十只股票。模型对每只股票返回一至三个分别描述的事件。每个事件取得五档方向强度之一，以及三档反应速度之一。若某只股票没有明确的方向性事件，模型返回一条中性记录。设股票 $i$ 的历史每小时波动率为 $\sigma_i$，则幅度为
\begin{equation}
 A_k=\ell_k\sigma_{i(k)},\quad
 \ell_k\in\{-2,-0.75,0,0.75,2\}.
\end{equation}
快速、中速和慢速反应的时间尺度分别为 $0.5$、$1.5$ 和 $4$ 小时。事件记录保留原因、来源轮次、类别和速度。每个事件独立贡献带符号的区间反应，因此同一只股票的预测可以同时包含方向相反的新闻。
相应的半衰时间约为 $0.31$、$0.92$ 和 $2.46$ 小时。离散等级和历史波动率限制了语言模型从简短摘要中指定收益幅度的范围。我们在开发日期上比较了幅度系数 $\{0.5,1,2\}$ 和三组时间尺度，见第五节；函数指数保持固定。

\subsection{股票关系偏离因子}
设 $\Sigma$ 为经过收缩估计的每小时协方差矩阵，$\rho_{ij}$ 为股票 $i$ 与 $j$ 的相关系数，$\beta_{ij}=\Sigma_{ij}/\Sigma_{jj}$。对于最近一次持续 $h_t$ 小时的当日对数收益向量 $r_t$，定义
\begin{align}
 d_{ij,t}&=r_{i,t}-\beta_{ij}r_{j,t},\\
 s_{ij,t}&=\sqrt{h_t(\Sigma_{ii}-\Sigma_{ij}\beta_{ij})}.
\end{align}
按照固定股票顺序，每个股票组合只枚举一次。当 $\rho_{ij}\geq0.45$ 且 $|d_{ij,t}|>2s_{ij,t}$ 时，该组合满足条件。因子的总幅度大小为
\begin{equation}
 B_{ij,t}=\rho_{ij}\min\{|d_{ij,t}|-2s_{ij,t},3s_{ij,t}\}.
\end{equation}
我们为股票 $i$ 追加幅度 $-\operatorname{sign}(d_{ij,t})B_{ij,t}/2$，为股票 $j$ 追加相反的幅度；时间尺度为 $\tau=1+4(1-\rho_{ij})$ 小时。关系因子仅使用当日交易时段内的观测值。
两个幅度之和为零。$2s$ 门槛排除较小偏离，$3s$ 上限限制极端观测值的影响，相关性越高，时间尺度越短。这些门槛属于建模假设；现有样本无法确定最优的关系衰减函数。

\subsection{因子汇总}
对每只股票，分别汇总正向贡献与负向贡献：
\begin{equation}
 a_{i,t}=\sum_{k:i(k)=i}\max(v_{k,t},0),\qquad
 b_{i,t}=\sum_{k:i(k)=i}\max(-v_{k,t},0).
\end{equation}
我们使用均匀收益分布 $U[-b_{i,t},a_{i,t}]$ 作为评分表示，由此得到
\begin{align}
 g_{i,t}&=\exp\{\mathbb{E}[\log(1+R_{i,t})]\}-1,\\
 \mu_{i,t}&=(a_{i,t}-b_{i,t})/2.
\end{align}
期望对数通过 16 点 Gauss--Legendre 求积计算，并要求 $b_{i,t}<1$。因子模块提供用于股票排序的 $g_{i,t}$，以及用于普通收益率计算的 $\mu_{i,t}$。

\subsection{根据结算后的预测进行反馈}
反馈规则将股票实际收益 $y_{i,t}$ 与决策时保存的因子模块几何预测 $g_{i,t}$ 比较：
\begin{equation}
 \epsilon_{i,t}=y_{i,t}-g_{i,t}.
\end{equation}
只有在 $|\epsilon_{i,t}|>2\sigma_i\sqrt{x_t-e_t}$ 时，股票才会交给语言模型解释；时间长度以小时计算。只要所需价格已经公布，这一规则适用于发生交易和保持空仓的决策所对应的预测。模型接收候选误差和当前摘要，判断是否追加因子。被接受的低、中、高三个强度分别对应 $\alpha\in\{0.25,0.5,1\}$。因子幅度为 $\alpha\epsilon_{i,t}$，时间尺度为三小时。因子方向由观察到的误差决定；模型决定误差是否需要在后续预测中持续体现，以及体现的强度。

三小时时间尺度对应的幂函数半衰时间约为 $1.76$ 小时。触发门槛同时考虑历史波动率与持仓时长，三个强度等级限制能够延续到后续预测中的误差比例。反馈使用原始因子预测、已经发布的价格、新闻和既有因子。因此，动量权重与交易门槛不会改变因子的生成过程。

\subsection{动量与对冲决策}
若某次发布在同一交易时段内存在前一次观测，则在 $N=10$ 只股票之间定义经过平均值调整的对数价格变化：
\begin{equation}
 m_{i,t}=\log\frac{P_{i,t}}{P_{i,t-1}}-
 \frac{1}{N}\sum_{j=1}^{N}\log\frac{P_{j,t}}{P_{j,t-1}}.
\end{equation}
若该发布是当日首次观测，则 $m_{i,t}=0$。决策规则按以下方式结合信号：
\begin{equation}
 q_{i,t}=g_{i,t}+\lambda m_{i,t},\qquad
 \widehat\mu_{i,t}=\mu_{i,t}+\lambda m_{i,t},\quad \lambda=1.5.
\end{equation}
系统选择 $q$ 最高的股票做多，选择 $q$ 最低的股票做空。只有两只股票不同，且满足
\begin{equation}
 \tfrac12(\widehat\mu_{L_t,t}-\widehat\mu_{S_t,t})-0.001>0.0005
\end{equation}
时才发生交易；其余情况保持空仓。令 $\lambda=0$ 可得到仅使用因子的交易参考方法。价格项使用最近已经观察到的变化；每个事件项表示未来进场至退出区间内尚存的反应。两类信息进入同一套确定性的股票选择和费用门槛计算。

\begin{table*}[!t]
\centering
\caption{在 37 个可以评价的轮次中，每次活跃的完整配对交易收取 10 个基点费用时的表现。收益和回撤单位为百分比；Sharpe 指标未经年化。}
\label{tab:performance}
\begin{tabular}{lrrrrrr}
\toprule
方法 & 开发日期 & 后续日期 & 全部日期 & 交易轮次 & Sharpe & 最大回撤\\
\midrule
TEMPO & 6.658 & 2.813 & 9.658 & 30 & 0.312 & 1.142\\
仅使用因子 & 4.103 & 0.548 & 4.673 & 7 & 0.305 & 0.149\\
FinBERT 加动量 & 2.140 & 2.957 & 5.160 & 37 & 0.200 & 2.284\\
\bottomrule
\end{tabular}
\end{table*}

\section{实验}
\subsection{评价方法}
10 月 14 日、15 日、19 日和 20 日构成四个日期的开发部分，含 22 个可以评价的轮次。10 月 21 日至 23 日构成后续日期部分，含 15 个可以评价的轮次。配置选择以开发部分的累计净收益最大为准。架构开发期间使用过这两个部分；分别报告其收益，用于描述研究样本内部随时间变化的表现。

因子参数校准包含九种仅使用事件的组合：将第 IV-B 节的强度等级分别乘以 $\{0.5,1,2\}$，并采用 $(0.5,1.5,4)$、$(1,3,8)$ 和 $(2,6,16)$ 小时三组时间尺度。反馈参数校准分别比较相对于几何预测、普通收益率预测，以及经过股票间平均值调整的普通收益率预测的误差。最终因子参数在第四节给出。

逐步扩展的交易决策搜索包含 15 组配置：四组仅使用因子的额外收益门槛，分别为 0、2.5、5 和 10 个基点；在 5 个基点门槛下设置五种动量权重，分别为 0.25、0.5、1、1.5 和 2；在权重为 1 时额外设置 0、2.5 和 10 个基点三种门槛；并在 5 个基点门槛下测试惩罚系数为 1、10 和 100 的三种在线岭回归配置。在线回归候选方法以经过股票间平均值调整的实际收益为目标，使用经过相同调整的普通收益率预测和价格变化作为输入；所有数值都除以 0.01，并通过正则化使系数趋近 $(1,0)$。它们仅使用退出价格已经公布的结果。权重 1.5 与 2 在开发部分取得相同的最高收益，并在可以评价的轮次中产生完全相同的交易行动。TEMPO 采用权重 1.5。

所有交易决策配置都使用同一份已经记录的因子生成过程。组成部分消融实验保留指定类型的因子，在动量权重 1.5、额外收益门槛 5 个基点的条件下重新计算预测和交易行动。不同方法使用相同的进场价格、退出价格和可以评价的轮次。累计收益定义为 $\prod_t(1+R_t^{\mathrm{net}})-1$。Sharpe 指标为所有已完成轮次收益的平均值除以其样本标准差，包含空仓轮次，且不进行年化。最大回撤按照结算时点计算，初始账户权益设为一。

\subsection{交易表现}
表~\ref{tab:performance} 将 TEMPO 与两种参考方法比较。仅使用因子的交易方法设定 $\lambda=0$，保留全部三类因子和 5 个基点的额外收益门槛。FinBERT 加动量的参考方法将 FinBERT 的正面情绪分数减去负面情绪分数\cite{araci2019}，再与最新的当日对数收益结合；该对数收益先减去股票间的中位数，再除以 0.01。该方法从检索到新闻的股票中选择两端的股票；若无法计算当日价格差，则只使用情绪分数。

TEMPO 的累计净收益为 9.66\%，比仅使用因子的交易方法高 4.98 个百分点。开发日期的收益为 6.66\%，后续日期的收益为 2.81\%。与仅使用因子的交易方法相比，交易轮次从七个增加至 30 个，按结算时点计算的最大回撤从 0.15\% 上升至 1.14\%。七个日期中有五个日期的收益得到改善。FinBERT 加动量的参考方法在全部日期取得 5.16\% 的收益，在后续日期取得 2.96\% 的收益；完整回放共交易 37 轮，按结算时点计算的最大回撤为 2.28\%。

\subsection{动量权重敏感性}
表~\ref{tab:execution} 在保持因子生成过程和 5 个基点的额外收益门槛不变时，改变动量权重。权重为 0.25 时，累计收益为 1.65\%；权重为 1.5 和 2 时，累计收益均为 9.66\%。权重 1、1.5 和 2 在后续日期取得相同收益；开发日期的股票选择使权重 1 与后两个权重产生差别。

\begin{table}[!tb]
\centering
\caption{额外收益门槛为 5 个基点、费用为 10 个基点时，结果对动量权重的敏感性。收益单位为百分比。}
\label{tab:execution}
\begin{tabular}{rrrrr}
\toprule
权重 & 开发日期 & 后续日期 & 全部日期 & 交易轮次\\
\midrule
0 & 4.103 & 0.548 & 4.673 & 7\\
0.25 & 1.183 & 0.463 & 1.652 & 20\\
0.5 & 4.955 & 0.458 & 5.436 & 26\\
1 & 6.223 & 2.813 & 9.210 & 30\\
1.5 & 6.658 & 2.813 & 9.658 & 30\\
2 & 6.658 & 2.813 & 9.658 & 30\\
\bottomrule
\end{tabular}
\end{table}

\subsection{价格与证据因子的贡献}
在动量权重 1.5 和固定费用门槛下，仅使用价格动量取得 3.15\% 的收益。加入事件因子后，收益提高至 9.40\%；再加入关系因子后，收益提高至 9.66\%（表~\ref{tab:components}）。四种组成部分配置均在相同的 30 个轮次交易，因此收益差异来自所选择的股票。在已经包含事件和关系因子时，反馈因子没有改变可以评价的轮次中的交易行动。事件和关系因子的收益贡献发生在开发日期部分；四种配置在后续日期的收益相同。

\begin{table}[!tb]
\centering
\caption{动量权重为 1.5、额外收益门槛为 5 个基点时，使用不同决策输入得到的结果。收益单位为百分比。}
\label{tab:components}
\begin{tabular}{lrrr}
\toprule
决策输入 & 开发日期 & 后续日期 & 全部日期\\
\midrule
价格 & 0.328 & 2.813 & 3.150\\
价格加事件 & 6.406 & 2.813 & 9.399\\
价格加事件、关系 & 6.658 & 2.813 & 9.658\\
价格加全部三类因子 & 6.658 & 2.813 & 9.658\\
\bottomrule
\end{tabular}
\end{table}

\subsection{费用与日期敏感性}
表~\ref{tab:cost} 在保持 TEMPO 的交易行动和决策门槛不变时，改变实际收取的费用。每个活跃轮次收取 10 个基点时，累计净收益为 9.66\%；收取 20 个基点时，累计净收益为 6.43\%。考察日期敏感性时，我们保持已经记录的反馈因子生成过程，每次从计分结果中移除一个日期。七种移除日期的情况中，剩余轮次的累计收益介于 6.57\% 与 9.85\% 之间。

\begin{table}[!tb]
\centering
\caption{保持交易行动不变时，TEMPO 对每次完整配对交易费用的敏感性。}
\label{tab:cost}
\begin{tabular}{rrr}
\toprule
费用（基点） & 累计收益（\%） & 最大回撤（\%）\\
\midrule
0 & 12.986 & 0.843\\
5 & 11.310 & 0.893\\
10 & 9.658 & 1.142\\
20 & 6.425 & 1.833\\
\bottomrule
\end{tabular}
\end{table}

\subsection{市场贝塔敞口}
我们使用雅虎财经的 SPY 基金与十只股票在相同时点的历史小时区间开盘价估计市场贝塔。样本包含 2023 年 12 月 19 日至 2026 年 10 月 13 日的 698 个完整交易日期和 3,490 个当日区间，均早于比赛首次发布。与协方差校准一致，对数收益除以区间小时数的平方根。普通最小二乘估计为 $\widehat\beta_i=\operatorname{Cov}(r_i,r_{\mathrm{SPY}})/\operatorname{Var}(r_{\mathrm{SPY}})$；等金额多空组合的估计贝塔为 $(\widehat\beta_L-\widehat\beta_S)/2$。历史回放中 30 次活跃且可评价的交易，其贝塔绝对值均值为 $0.488$，中位数为 $0.479$，最大值为 $0.907$。五次实时活跃交易的贝塔绝对值均值为 $0.568$，最大值为 $0.907$。因此，等金额配置仍存在明显的估计市场敞口。SPY 只是市场代理，历史估计无法直接衡量持仓区间内实际出现的敞口。

\section{实现与可复现性}
实现使用 DeepSeek API 中标识为 \texttt{deepseek-flash} 的模型，关闭 thinking，并要求以 JSON 对象返回结果。事件和反馈请求的输出上限分别为 3,200 与 1,600 个 token。提示词要求模型仅使用提供的数据。输出必须满足股票顺序、类别、数量与类型约束；无效响应会导致执行立即终止。回放进行了 41 次事件请求和八次反馈请求，共使用 36,940 个输入 token 和 15,410 个输出 token。

请求缓存的键由模型标识和完整消息的 SHA-256 摘要构成。缓存记录响应标识、返回的模型标识、fingerprint、解析后的输出以及 token 用量。报告中的交易决策比较重现同一份已保存的模型响应记录。每份报告都记录输入文件哈希、数值参数、因子、预测、交易行动与结算结果。

独立回放使用已经公布的进场和退出价格，重新构建因子贡献、股票选择与收益。在每次发布时，根据完整已保存价格数组算出的决策，与将价格数组截断至该发布时点后算出的决策相同。另一项在第 25 轮结束的智能体运行，重现了 41 轮运行中前 25 轮的决策、因子和预测。反馈资格按照每个退出价格观测值的发布时间检查。这些检查验证了报告中的数值结果与时间访问规则。

\FloatBarrier
\section{实时交易}
\subsection{前瞻性规则与固定配置}
第三阶段于 2026 年 11 月 2 日进行六轮独立计分的交易，初始账户权益重新设为一\cite{task2protocol,task2livefinal}。预定发布时刻为纽约时间 09:30 至 14:30，每小时一次，时区为 EST（UTC$-5$）。每个行动在下一个估值时点进场，最终持仓区间为 15:30--16:00。官方最终数据在 EST 16:03:14 更新，确认全部六轮已经核验，阶段状态为结束。

我们保留交易开始前确定的 TEMPO 配置：使用 \texttt{deepseek-flash}，关闭 thinking，沿用第四节的因子参数，动量权重为 1.5，决策费用预留为 10 个基点，额外收益门槛为 5 个基点。保存的方法哈希与当前实现一致，六份决策记录使用相同模型和交易规则。因子历史从 10 月回放延续，11 月 2 日的成绩单独计算。随着数据发布，程序追加新的因子与符合条件的反馈。官方行动历史与保存的六次决策一致：第一轮空仓，随后五轮执行做多、做空组合（表~\ref{tab:live_actions}）。

\begin{table}[!tb]
\centering
\caption{TEMPO 实时行动与可取得的官方逐轮收益。时间为纽约 EST；股票组合依次表示做多与做空。横线表示该轮单独收益不可取得。}
\label{tab:live_actions}
\small
\setlength{\tabcolsep}{3pt}
\begin{tabular}{cllr}
\toprule
发布时刻 & 持仓区间 & 交易行动 & 收益（\%）\\
\midrule
09:30 & 10:30--11:30 & 空仓 & 0.000\\
10:30 & 11:30--12:30 & AMZN/AAPL & 0.446\\
11:30 & 12:30--13:30 & WMT/AMD & 0.826\\
12:30 & 13:30--14:30 & WMT/AAPL & 0.162\\
13:30 & 14:30--15:30 & GOOG/AMD & ---\\
14:30 & 15:30--16:00 & WMT/AMD & ---\\
\bottomrule
\end{tabular}
\end{table}

\subsection{官方成绩与证据范围}
表~\ref{tab:live_results} 列出 TEMPO 与三种官方基线的最终指标\cite{task2livefinal}。TEMPO 的累计收益为 2.274751\%，非年化 Sharpe 为 1.294271，结算时点最大回撤为零。TEMPO 在参赛队伍中排名第二，累计收益超过 Moira、FinMem 和 Codex。

\begin{table}[!tb]
\centering
\caption{TEMPO 与官方基线的六轮实时结果。收益与结算时点最大回撤的单位为百分比；Sharpe 未年化。}
\label{tab:live_results}
\begin{tabular}{lrrr}
\toprule
队伍或基线 & 收益 & Sharpe & 最大回撤\\
\midrule
TEMPO & 2.275 & 1.294 & 0.000\\
Moira（基线） & 1.063 & 0.458 & 0.094\\
FinMem（基线） & 0.669 & 0.715 & 0.166\\
Codex（基线） & $-1.517$ & $-0.668$ & 1.707\\
\bottomrule
\end{tabular}
\end{table}

官方指标只使用六轮实时交易结果。Sharpe 包含空仓轮次，采用样本标准差，无风险收益设为零；最大回撤包含初始权益和各结算时点\cite{task2protocol}。这里直接报告组织方成绩，研究回放和决策门槛中的 10 个基点费用没有再次扣除。本地报告跨越 10 月与 11 月的累计结果不计入本节实时评价。

前四轮的收益可以从连续保存的官方快照计算，依次为 0.000000\%、0.446289\%、0.825831\% 和 0.162112\%，与根据公开估值价格独立计算的等金额多空毛收益一致。最终快照提供汇总指标与六轮行动，但没有第 5--6 轮各自的精确收益。将最终权益除以第四轮结束时的权益，再减去一，可以得到这两轮合计复利收益 0.822916\%。保存的公开价格发布截至 14:30，因此无法从这些文件独立检查最终 15:30 和 16:00 价格。本文保留官方最终指标，并将这两轮各自的收益标为不可取得。

\subsection{结果解释与局限}
本次交易使用数据发布前已经固定的方法，记录了遵守时间访问规则的执行过程。五次活跃交易覆盖单个交易日和十只股票，样本不足以确立持续盈利能力、统计显著的优势，或单独某类因子在实时交易中的贡献。等金额多空组合仍具有非零的估计市场贝塔；零结算时点回撤也无法衡量持仓区间内部的风险。比赛采用模拟估值和执行，实际证券账户的成交、滑点与融资费用不在本次结果的评价范围内。评价泛化能力和预测校准还需要更长的前瞻性观察，以及完整的结算数据。

\FloatBarrier
\bibliographystyle{IEEEtran}
\section*{代码与数据}
实现程序、提示词、缓存的模型解析结果、公开输入及数值报告见 \url{https://github.com/shenma32/TEMPO}。

\section*{致谢}
OpenAI Codex 辅助起草和修订了摘要与正文（第一至第七节）。作者对最终内容负责。

\bibliography{official-overviews,official-herculean,references,xie-finance}

\clearpage
\section*{面向初学者的完整解读}

\subsection*{一、比赛要求完成什么}
比赛会分批公布十只股票的价格和简短新闻。每次看到新的公开内容，程序都要决定是否交易；若交易，就指定一只股票做多、另一只股票做空。程序作出决定的时刻与真正进场的时刻之间有间隔。进场以后，还要等到预定的下一个时点退出。因此，看到一条新闻时，关键问题是：等到实际持仓开始，这条新闻对价格可能还有多少影响？

做多可以理解为希望股票上涨；做空可以理解为希望股票下跌。本文的方法把一半账户权益分配给做多方向，另一半分配给做空方向。例如，做多股票在持仓期间上涨 2\%，做空股票在持仓期间上涨 1\%，交易的毛收益为 $(2\%-1\%)/2=0.5\%$。如果做空股票下跌 1\%，毛收益则为 $(2\%-(-1\%))/2=1.5\%$。这种组合仍然会受到两只股票各自价格变化的影响，因此“市场中性”表达的是对冲设计，并不保证收益稳定。

\subsection*{二、程序在每次发布时做什么}
程序先检查此前的预测是否已经有可以使用的结算结果。符合条件的较大预测误差可能产生反馈因子。接着，程序从本次新闻提取事件因子，从相关股票的价格变化提取关系因子。因子可以理解为一条记录：它针对哪只股票、倾向上涨还是下跌、影响有多强、从什么时候开始，以及影响消退得有多快。记录创建后保留原样，新证据通过追加记录进入系统。

程序再计算各条记录在未来持仓期间预计发挥的作用，并结合最新的相对价格变化，为十只股票打分。最后选出分数最高的股票做多、分数最低的股票做空。只有预计收益在扣除交易费用后还超过额外门槛，程序才执行交易；其余轮次保持空仓。

\subsection*{三、为什么新闻需要考虑时间}
一条利好消息发布以后，部分价格反应可能在程序真正进场以前就已经发生。本文用一个随时间下降的函数表示“尚未发生的影响”。在进场时刻查看剩余多少，在退出时刻再查看剩余多少，两者的差额就是这段持仓期内预计发生的部分。事件、股票关系偏离和反馈分别采用不同的衰减函数。时间按照实际经过的小时计算，因此夜间也计入消息变旧的过程。

举一个只用于解释公式的数值例子：某条利好消息的幅度是 1\%，进场时尚存影响比例为 60\%，退出时为 40\%，则该消息在持仓期间的预测贡献为 $1\%\times(60\%-40\%)=0.2\%$。实际计算使用正文给出的三个函数和每条记录的时间尺度；这个例子没有对应比赛中的具体交易轮次。

时间尺度 $\tau$ 决定影响消退的快慢。新闻事件使用 0.5、1.5 或 4 小时；关系因子的时间尺度随相关系数变化；反馈因子使用三小时。公式中的正负号表示预期方向，尚存影响比例随着时间降低，所以区间贡献沿用原记录的方向。

\subsection*{四、三类因子分别表达什么}
新闻事件因子来自语言模型对摘要的解释。模型为每个事件指定方向强度和反应速度；历史每小时波动率决定强度对应的数值幅度。同一只股票若同时出现利好和利空，两条记录可以同时存在。语言模型负责解释文字，后续数值计算遵循固定公式。

关系因子考察两只通常有相关价格变化的股票。程序先估计它们的历史协方差，再查看最近的当日价格变化。如果其中一只股票相对另一只股票的变化超过统计阈值，就记录一个方向相反的股票组合信号。这个信号表达的是模型对异常偏离未来变化的判断，价格仍然可能继续偏离。

协方差表示两只股票的价格变化是否经常一起出现，相关系数把这种关系缩放到便于比较的数值范围。正文中的 $\beta_{ij}$ 用历史协方差估计股票 $j$ 变化时股票 $i$ 通常会怎样变化。$d_{ij,t}$ 是股票 $i$ 的实际变化减去据此估计的变化；$s_{ij,t}$ 是这种差值通常波动的大小。只有相关系数至少为 0.45，且偏离超过通常波动大小的两倍，程序才建立关系因子。因子幅度还受到通常波动大小的三倍这一上限约束。

反馈因子来自已经结算的预测。程序保存当时的预测，等到退出价格正式发布后再计算实际收益与预测的差距。只有差距超过由历史波动率决定的阈值，才交给语言模型判断是否值得保留为后续信号。保持空仓的轮次也会产生可检查的股票预测，只要对应价格已经公布。

把所有因子汇总时，程序先分别累加预计上涨和预计下跌的贡献，得到 $a_{i,t}$ 与 $b_{i,t}$。然后，它用从 $-b_{i,t}$ 到 $a_{i,t}$ 的均匀分布表示这只股票可能取得的收益。该分布给出两种预测：$\mu_{i,t}$ 是普通收益率的平均值，$g_{i,t}$ 则先平均对数收益，再换算成收益率。程序用 $g_{i,t}$ 参与股票排序，用 $\mu_{i,t}$ 检查配对收益。正文提到的 16 点 Gauss--Legendre 求积，是计算上述对数平均值的一种数值积分方法；这里的“16 点”表示计算时选用的 16 个积分位置。

\subsection*{五、价格动量与交易门槛怎样发挥作用}
价格动量比较同一交易时段内十只股票最近的价格变化。程序先计算各股票的对数价格变化，再减去十只股票的平均变化。结果为正，表示这只股票近期表现高于这组股票的平均水平；结果为负，则表示低于平均水平。当天第一次观测时，没有可用的当日上一时点，动量设为零。

最终排序分数等于因子预测加上 $1.5$ 倍的相对价格动量。程序还用相应的普通收益率预测检查交易门槛。一个基点等于 0.01\%：10 个基点的费用等于账户权益的 0.10\%，额外的 5 个基点门槛等于 0.05\%。因此，正文中的条件要求预测的配对毛收益超过 0.15\%。这个门槛只控制是否交易，并不保证真实收益超过费用。

正文采用对数价格变化，因为连续价格变化可以方便地相加。读者可以把这里的数值理解为较短时间内的价格变化比例；程序始终按照正文给出的对数公式计算。动量权重 1.5 决定近期相对价格变化在评分中的影响大小。表二展示其他权重的回放结果，权重改变后，股票选择和交易轮次都可能改变。

\subsection*{六、收益数字怎样得到}
本次回放含 41 个决策轮次；其中 37 个轮次具有完整的进场和退出价格，可以计算收益。其余四个轮次缺少最终价格，因此没有计入收益。TEMPO 在可以评价的轮次中交易 30 次。每次活跃交易按账户权益收取 0.10\% 费用，空仓轮次不收取费用。把各轮净收益依次作用于账户权益后，全部日期的累计净收益为 9.658\%，正文四舍五入写作 9.66\%。这表示研究样本内按规定价格回放得到的结果。

只使用因子时，累计净收益为 4.673\%，交易七次。只使用价格动量时，累计净收益为 3.150\%；加入新闻事件后为 9.399\%；再加入关系因子后为 9.658\%。在这一批可评价轮次中，加入反馈因子没有改变交易行动。由于这些方法都在相同的 30 个轮次执行交易，组成部分实验中的收益差异来自股票选择。将每次交易费用提高到 0.20\%，保持原有交易行动，累计净收益变为 6.425\%。

\subsection*{七、怎样理解风险和实验范围}
最大回撤衡量账户权益从此前最高点到随后较低点的最大下降幅度。本文报告的 1.142\% 只根据各轮结算时点计算，无法描述两次结算之间可能出现的价格波动。Sharpe 指标使用每个已完成轮次的净收益，包括空仓轮次，且没有转换成年化数值。

正文的“消融实验”指每次保留指定的决策输入，重新计算股票选择和收益，由此观察各类信息在这批数据中的贡献。“费用敏感性”则固定已经选好的交易行动，只改变结算时扣除的费用。表格中的开发日期收益、后续日期收益和全部日期收益，分别对应不同的日期范围；全部日期的累计收益需要逐轮复利计算，直接相加两个日期范围的百分比会产生少量误差。

10 月 14 日、15 日、19 日和 20 日的 22 个可评价轮次用于配置选择；10 月 21 日至 23 日的 15 个可评价轮次单独报告。架构研究期间，两个日期部分都曾被使用。因此，后续日期结果可以帮助阅读者观察样本内的时间变化，不能直接当作完全独立的未来检验。回放也无法直接证明真实交易会取得同样的成交价格或收益。

\subsection*{八、怎样核查程序有没有使用未来信息}
历史参数所用数据截止于比赛首次发布之前。每一轮只允许使用已经发布的新闻和价格；预测误差必须等相应退出价格公布后才能进入反馈。独立回放重新计算因子、股票选择与收益；把价格数据截断至当前发布时点后，计算出的决策与使用完整已保存数组时相同。在第 25 轮停止的单独运行，也重现了完整运行的前 25 轮决策、因子和预测。这些检查针对数据使用时间和数值计算的一致性。

\subsection*{九、实时成绩与研究回放怎样分别阅读}
11 月 2 日的六轮交易在新数据逐次发布时作出决定，使用事前固定的模型与参数。第一轮保持空仓，后续依次做多 AMZN、WMT、WMT、GOOG、WMT，并分别做空 AAPL、AMD、AAPL、AMD、AMD。官方最终行动与本地保存记录一致。第三阶段以账户权益一重新开始，最终权益为 1.022747512，对应官方累计收益 2.274751\%。这项成绩只覆盖当天六轮；10 月研究回放的 9.658\% 单独报告。

TEMPO 在参赛队伍中排名第二，收益高于全部三种官方基线。官方 Sharpe 为 1.294271，表示六轮收益的平均值与样本标准差之比，空仓轮次也参与计算；它没有年化。最大回撤为零，表示结算时点的权益序列没有从此前高点下降，持仓期间仍然可能发生价格波动。

前四轮收益可以使用当时保存的官方累计成绩逐轮还原，也能通过公开价格重新计算。最后两轮的合计复利收益为 0.822916\%，但公开最终快照没有给出两轮各自的精确收益，保存的价格文件也缺少最后两个估值时点。正文因此明确记录可取得的数据范围。单个交易日的五次活跃交易提供了实际按时决策的证据，对长期表现与各因子的独立贡献仍需更长观察。

\end{document}
